\section{Problem Analysis}

The problem presents several key challenges. The model must capture both long-range temporal dependencies and cross-channel interactions, while integrating heterogeneous features, including survey, biomechanical, and global signals with different statistical properties. The limited dataset size increases the risk of overfitting, requiring regularization, careful hyperparameter tuning, and robust cross-validation. Additionally, class imbalance may bias predictions toward frequent labels, which we mitigate using inverse-frequency class weights and stratified folds.

Our approach is based on several assumptions. Pain levels are primarily reflected in temporal dynamics, including trends, variability, and motion irregularities, which may occur at multiple temporal scales from slow posture drifts to fast joint oscillations, motivating the use of convolutional and multiscale encoders. Global static features, such as subject-specific metadata, provide priors that influence pain prediction and should be integrated with dynamic signals. Dimensionality reduction techniques, like PCA or POD, can capture dominant modes, reduce redundancy, and improve generalization. Finally, given the dataset size, K-fold cross-validation is essential for reliable performance estimates. These considerations guided the design of our modular preprocessing pipeline and flexible architecture search.

